% Options for packages loaded elsewhere
\PassOptionsToPackage{unicode}{hyperref}
\PassOptionsToPackage{hyphens}{url}
\documentclass[
  11pt,
]{article}
\usepackage{xcolor}
\usepackage[margin=2.2cm]{geometry}
\usepackage{amsmath,amssymb}
\setcounter{secnumdepth}{-\maxdimen} % remove section numbering
\usepackage{iftex}
\ifPDFTeX
  \usepackage[T1]{fontenc}
  \usepackage[utf8]{inputenc}
  \usepackage{textcomp} % provide euro and other symbols
\else % if luatex or xetex
  \usepackage{unicode-math} % this also loads fontspec
  \defaultfontfeatures{Scale=MatchLowercase}
  \defaultfontfeatures[\rmfamily]{Ligatures=TeX,Scale=1}
\fi
\usepackage{lmodern}
\ifPDFTeX\else
  % xetex/luatex font selection
  \setmainfont[]{DejaVu Serif}
  \setsansfont[]{DejaVu Sans}
  \setmonofont[]{DejaVu Sans Mono}
\fi
% Use upquote if available, for straight quotes in verbatim environments
\IfFileExists{upquote.sty}{\usepackage{upquote}}{}
\IfFileExists{microtype.sty}{% use microtype if available
  \usepackage[]{microtype}
  \UseMicrotypeSet[protrusion]{basicmath} % disable protrusion for tt fonts
}{}
\makeatletter
\@ifundefined{KOMAClassName}{% if non-KOMA class
  \IfFileExists{parskip.sty}{%
    \usepackage{parskip}
  }{% else
    \setlength{\parindent}{0pt}
    \setlength{\parskip}{6pt plus 2pt minus 1pt}}
}{% if KOMA class
  \KOMAoptions{parskip=half}}
\makeatother
\setlength{\emergencystretch}{3em} % prevent overfull lines
\providecommand{\tightlist}{%
  \setlength{\itemsep}{0pt}\setlength{\parskip}{0pt}}
% ---- custom preamble for the RoadFury -> SE2RoadNet speaker script ----
\usepackage{xcolor}
\usepackage{framed}
\usepackage{titlesec}
\usepackage{enumitem}
\usepackage{parskip}

\definecolor{deckteal}{HTML}{16242A}
\definecolor{deckorange}{HTML}{E0772B}
\definecolor{deckslate}{HTML}{45595B}

\setcounter{secnumdepth}{0}
\linespread{1.05}

% coloured, unnumbered headings matching the deck palette
\titleformat{\section}{\normalfont\Large\bfseries\color{deckorange}}{}{0pt}{}
\titleformat{\subsection}{\normalfont\large\bfseries\color{deckteal}}{}{0pt}{}
\titlespacing*{\section}{0pt}{1.5em}{0.55em}
\titlespacing*{\subsection}{0pt}{1.05em}{0.35em}

% markdown blockquotes ( "> VERIFY: ..." ) -> small italic left-bar callout
\renewenvironment{quote}
  {\par\smallskip\noindent
   \def\FrameCommand{{\color{deckorange}\vrule width 2.5pt}\hspace{10pt}}%
   \MakeFramed{\advance\hsize-\width\FrameRestore}%
   \footnotesize\itshape\color{deckslate}}
  {\par\endMakeFramed\smallskip}

\setlist{topsep=2pt,parsep=1pt,itemsep=1pt,leftmargin=1.4em}

% DejaVu Serif lacks U+2713/U+2717; borrow the check/cross from DejaVu Sans
\usepackage{newunicodechar}
\newfontfamily\dvsans{DejaVu Sans}
\newunicodechar{✓}{{\dvsans ✓}}
\newunicodechar{✗}{{\dvsans ✗}}
\usepackage{bookmark}
\IfFileExists{xurl.sty}{\usepackage{xurl}}{} % add URL line breaks if available
\urlstyle{same}
% fallback for those not using the hyperref driver hyperxmp:
\makeatletter
\@ifundefined{xmpquote}{\newcommand{\xmpquote}[1]{#1}}{}
\makeatother
\hypersetup{
  pdftitle={RoadFury → SE2RoadNet --- Script thuyết trình},
  pdfauthor={Trần Chí Nguyên · Huỳnh Trung Kiệt · Đào Sỹ Duy Minh},
  hidelinks,
  pdfcreator={LaTeX via pandoc}}

\title{RoadFury → SE2RoadNet --- Script thuyết trình}
\usepackage{etoolbox}
\makeatletter
\providecommand{\subtitle}[1]{% add subtitle to \maketitle
  \apptocmd{\@title}{\par {\large #1 \par}}{}{}
}
\makeatother
\subtitle{The Self-Driving Car Testing Competition · ICST 2026 · deck 37
slide (s00--s36)}
\author{Trần Chí Nguyên · Huỳnh Trung Kiệt · Đào Sỹ Duy Minh}
\date{Nguồn: các section HTML trong slide/google-slides · biên dịch
XeLaTeX}

\begin{document}
\maketitle

{
\setcounter{tocdepth}{1}
\tableofcontents
}
\begin{quote}
Deck: \texttt{presentation/slide/RoadFury-to-SE2RoadNet.pptx} (s00--s31
trình bày + s32--s36 backup cho Q\&A). Giới hạn của thầy: \textbf{tối đa
15 phút trình bày + 10 phút Q\&A}, trọng tâm \textbf{Method (§02--03)}
và \textbf{Benchmark (§04)}. Ngân sách thời gian gợi ý bên dưới: phần
Bối cảnh (§01) đi nhanh để dồn thời gian cho Method + Benchmark.

Mỗi mục: \textbf{{[}thời lượng{]}} + lời nói (đọc tự nhiên, không đọc
nguyên si) + dòng \texttt{VERIFY:} nếu trang đó có số liệu cần lưu ý.
\end{quote}

\begin{center}\rule{0.5\linewidth}{0.5pt}\end{center}

\subsection{(!) TRƯỚC KHI THUYẾT TRÌNH --- 5 điểm kiểm chứng phải xử
lý}\label{trux1b0ux1edbc-khi-thuyux1ebft-truxecnh-5-ux111iux1ec3m-kiux1ec3m-chux1ee9ng-phux1ea3i-xux1eed-luxfd}

Đối chiếu deck với dữ liệu gốc trong repo, có 5 chỗ \textbf{không khớp /
không có nguồn}. Đọc kỹ trước, vì thầy sẽ hỏi đúng những chỗ này.

\begin{enumerate}
\def\labelenumi{\arabic{enumi}.}
\item
  \textbf{``Cắt 50--80\% thời gian/chi phí'' (s03, s05, s35)} ---
  \textbf{{[}OK{]} ĐÃ SỬA (2026-06-24).} Đã thay bằng phát biểu định
  tính suy ra được từ chính APFD của nhóm (không cần cite): ``dồn lỗi
  lên đầu hàng đợi → chỉ cần chạy phần đầu đã lộ phần lớn lỗi → tiết
  kiệm chi phí mô phỏng''. Lý do bỏ con số: grep toàn bộ \texttt{exps/}
  = 0 kết quả, ``50--80\%'' không phải số nhóm đo. (Nếu muốn con số cụ
  thể của nhóm: Exp 05 cho APFD@K=50/287 ≈ 0.94 → chạy
  \textasciitilde17\% test đầu thu gần hết giá trị phát hiện lỗi.)
\item
  \textbf{Leaderboard RoadFury 0.804 \textless{} SE2RoadNet 0.805 (s07,
  s09, s23, s24)} --- \textbf{ĐẢO NGƯỢC dữ liệu gốc.} Số ``to beat''
  chính thức của dự án là RoadFury / best-single baseline =
  \textbf{0.8066 ± 0.0124}; SE2RoadNet = \textbf{0.8048 ± 0.0118}
  (\texttt{tracker.md} §3 Exp 02: ``APFD-comp 0.8048 is 0.0018 BELOW
  best-single baseline 0.8066''). Tức trên APFD, SE2RoadNet \textbf{thấp
  hơn} RoadFury 0.0018 --- nằm trong 1σ (≈0.012), nên đúng là
  \textbf{hòa trong sai số}, KHÔNG phải thắng.

  \begin{itemize}
  \tightlist
  \item
    Số 0.8066 làm tròn ra \textbf{0.807}, không phải 0.804. Con số 0.804
    chỉ ứng với cấu hình γ=1.5 (0.8045), không phải best-single γ=2.5
    (0.8066).
  \item
    \textbf{Sửa cho trung thực:} để RoadFury ≈ 0.807 (hoặc ghi rõ ``hòa
    trong σ''), SE2RoadNet 0.805. Giữ headline thật = \textbf{AUC ↑
    (0.917→0.934) + Δ=0} = Pareto improvement. (Script Q\&A cũ Q3 đã nói
    đúng tinh thần này: ``Em không claim APFD tốt hơn'' --- phần hình
    ảnh phải khớp với câu đó.) → Đây là điểm dễ bị bắt nhất. Đừng để bar
    chart ngụ ý SE2RoadNet thắng APFD.
  \end{itemize}
\item
  \textbf{Bảng xoay ITS4SDC (s13/s23):
  0.7810→0.7240/0.7518/0.7396/0.7334/0.7627, Δ=0.057} ---
  \texttt{tracker.md} Exp 02 chỉ ghi baseline ``should drop 0.04--0.07''
  như \textbf{dự đoán} + TODO ``Run the baseline rotation probe''; không
  tìm thấy bảng 6 góc đã đo. \textbf{THÊM:} paper ITS4SDC định nghĩa góc
  là ``góc giữa 2 đoạn kề'' (turning angle, \textbf{vốn bất biến xoay})
  → bản đúng-paper sẽ KHÔNG drift. Δ=0.057 chỉ đúng nếu bản nhóm tái lập
  dùng \textbf{góc/heading tuyệt đối} (phụ thuộc khung). → Phải có log
  tái lập sẵn cho Q\&A; nếu thầy biết ITS4SDC, đây là chỗ bị hỏi.
\item
  \textbf{Số baseline --- ĐÃ TRA LITERATURE (2026-06-24):}

  \begin{itemize}
  \tightlist
  \item
    {[}OK{]} \textbf{Greedy-diversity 0.795} --- KHỚP CHÍNH XÁC.
    Birchler et al., TOSEM 2023 ({[}2{]}): greedy search average APFD =
    \textbf{79.5\%}.
  \item
    {[}OK{]} \textbf{SO-SDC-Prioritizer 0.765} --- KHỚP CHÍNH XÁC. TOSEM
    2023: SO-SDC average APFD = \textbf{76.5\%}.
  \item
    (!) \textbf{Random 0.493} --- literature ghi \textbf{49.9\%} (TOSEM
    2023). Lệch nhẹ (0.493 vs 0.499); 0.493 là số nhóm tái lập, nói ``≈
    random \textasciitilde0.5'' là an toàn.
  \item
    {[}OK{]} \textbf{ITS4SDC (đã sửa tên 2026-06-24, trước ghi nhầm
    ``ITS4SDC'').} Güllü, Shah, Pfahl --- ICST 2025 SDC Tool Competition
    (arXiv:2501.03881). Là \textbf{bi-LSTM} (220 cells) trên 2 đặc trưng
    chuỗi (góc đoạn + chiều dài), test SELECTION, paper công bố
    \textbf{F1=0.89} --- \textbf{KHÔNG có APFD}. → APFD 0.781 + bảng
    xoay Δ=0.057 là \textbf{nhóm tái lập}, không phải số trong paper.
    Slide đã sửa: tên ITS4SDC + method LSTM (bỏ ``MLP 3 đặc trưng'').
    Còn lại: drift Δ=0.057 cần log tái lập (xem (!) 3) + sửa nhãn
    ``ICST'25'' → ref riêng cho ITS4SDC (hiện refs deck chưa có entry
    ITS4SDC --- nên thêm).
  \item
    (!) \textbf{GNN 0.533, ResNet-50 0.572, LLM zero-shot 0.487} ---
    không có nguồn công bố; là baseline nhóm tái lập trong harness. Cần
    sẵn log để chứng minh.
  \end{itemize}
\item
  \textbf{Resolution probe Δ=0.0012 (s15 caption, s27, s34/B3)} --- số
  này đo trên \textbf{mô hình FNO (Exp 01)}, KHÔNG phải SE2RoadNet,
  nhưng deck SE2RoadNet trình bày như thể của SE2RoadNet. Ngoài ra bar ở
  B3 (.8060/.8066/.8072/.8067) \textbf{không khớp} bảng FNO thật
  (.8051/.8060/.8063/.8060/.8062) --- cùng ra range 0.0012 nhưng giá trị
  từng N khác. → Hoặc chạy probe trên SE2RoadNet, hoặc ghi rõ ``FNO (Exp
  01)'', và dùng đúng số bảng thật.
\end{enumerate}

\textbf{Số đã kiểm và ĐÚNG (cứ tự tin nói):} dataset 28804/7202/956 +
38.4/38.4/ 36.9\% ✓; SE2RoadNet 2.11M params (2,108,721) ✓,
\textasciitilde2.5× RoadFury ✓, train 24.2 phút ✓, Focal γ=1.5 ✓,
d=192/6 lớp/8 heads/32 RFF ✓; RoadFury 829K params (828,801) ✓, AUC
0.917 ✓; multi-trial 30 lần, mẫu 287/956 = max(50,
0.3·\textbar test\textbar) ✓, seed 42 ✓, SE2RoadNet 0.8048±0.0118 ✓;
single-pass 0.8047 ✓; rotation Δ=0.0000/6 góc ✓; AUC 0.9347 cao nhất dự
án ✓ (lưu ý làm tròn ra \textbf{0.935}, deck ghi 0.934 --- sai số làm
tròn, nên dùng 0.935 hoặc 0.9347); công thức APFD + ví dụ tay B1 (0.80 /
0.20) ✓; AUC/APFD phân kỳ qua 4 exp ✓; toàn bộ trang Hạn chế (s27) khớp
tracker ✓.

\begin{center}\rule{0.5\linewidth}{0.5pt}\end{center}

\section{PHẦN MỞ ĐẦU}\label{phux1ea7n-mux1edf-ux111ux1ea7u}

\subsection{s00 · Title --- {[}15s{]}}\label{s00-title-15s}

Em chào thầy và cả lớp. Nhóm em trình bày đề tài \textbf{ưu tiên kiểm
thử xe tự lái} (self-driving car test prioritization). Tên phương pháp
là \textbf{RoadFury → SE2RoadNet}: RoadFury là baseline nền tảng của
nhóm, SE2RoadNet là phiên bản cải tiến có \textbf{bất biến hình học
SE(2)} --- xoay hay dịch con đường thì điểm số không đổi. Nhóm gồm 3
thành viên.

\subsection{s01 · Outline --- {[}20s{]}}\label{s01-outline-20s}

Bài có 5 phần: (1) bối cảnh và bài toán, (2) RoadFury --- phương pháp
nền, (3) SE2RoadNet --- đề xuất chính, (4) đánh giá thực nghiệm, (5) hạn
chế và hướng phát triển. Trọng tâm là phần \textbf{3 và 4}. Em có thêm 5
slide backup cho phần hỏi đáp.

\begin{center}\rule{0.5\linewidth}{0.5pt}\end{center}

\section{§01 --- BÀI TOÁN (đi nhanh, \textasciitilde2.5 phút
tổng)}\label{buxe0i-touxe1n-ux111i-nhanh-2.5-phuxfat-tux1ed5ng}

\subsection{s02 · Divider §01 --- {[}5s{]}}\label{s02-divider-01-5s}

Phần một: bài toán.

\subsection{s03 · Bối cảnh ---
{[}40s{]}}\label{s03-bux1ed1i-cux1ea3nh-40s}

Quy trình test xe tự lái hiện nay chạy trong \textbf{mô phỏng} như
BeamNG.tech: mỗi build sinh ra hàng nghìn kịch bản đường, mỗi kịch bản
chạy vật lý mất nhiều giây tới một phút, cộng lại là \textbf{nhiều giờ
CPU}. Phần lớn kịch bản đều PASS, chỉ \textasciitilde30--40\% là FAIL.
Nên thay vì chạy ngẫu nhiên, ta \textbf{xếp các kịch bản dễ FAIL lên
đầu} để lộ lỗi sớm.

\begin{quote}
\textbf{VERIFY:} ``cắt 50-80\% thời gian'' ở cuối card 3 ---
\textbf{không có nguồn} (xem điểm (!) 1). Nói ``cắt phần lớn chi phí khi
chỉ chạy prefix ngắn'' cho an toàn, hoặc dẫn số APFD@K của ta.
\end{quote}

\subsection{s04 · Phát biểu bài toán ---
{[}40s{]}}\label{s04-phuxe1t-biux1ec3u-buxe0i-touxe1n-40s}

Hình thức hóa: \textbf{đầu vào} là một con đường = chuỗi N điểm 2D (N từ
64 đến 197), chỉ có hình học, không có hành vi xe. \textbf{Scorer} f\_θ
ánh xạ con đường → xác suất FAIL trong {[}0,1{]}. \textbf{Đầu ra} là
hoán vị π xếp giảm dần theo điểm, mục tiêu tối đa hóa \textbf{APFD}. Khó
ở chỗ: chỉ có hình học đường, và tập thử là out-of-distribution.

\begin{quote}
\textbf{VERIFY:} N 64--197 ✓ khớp dataset. Metric APFD (chính) + AUC
(phụ) ✓.
\end{quote}

\subsection{s05 · Vì sao cần ưu tiên ---
{[}35s{]}}\label{s05-vuxec-sao-cux1ea7n-ux1b0u-tiuxean-35s}

APFD chính là \textbf{diện tích dưới đường cong ``đã phát hiện bao nhiêu
lỗi sau khi chạy x\% test''}. Ranking tốt → lỗi dồn lên đầu → đường cong
dốc đứng sớm. Ranking kém → lỗi nằm cuối → tốn nhiều giờ mô phỏng. Với
956 kịch bản nhân mô phỏng vật lý, chạy hết là rất đắt --- nên bài toán
không phải ``chạy nhanh hơn'' mà là ``chạy đúng thứ tự''.

\begin{quote}
\textbf{VERIFY:} takeaway ``cắt 50-80\% chi phí'' --- \textbf{cùng vấn
đề (!) 1}, xử lý như s03. ``APFD ngẫu nhiên \textasciitilde0.52'' trong
chú thích: random thật \textasciitilde0.49--0.50; 0.52 hơi cao nhưng chỉ
minh họa, không sai bản chất.
\end{quote}

\subsection{s06 · Dataset SensoDat ---
{[}45s{]}}\label{s06-dataset-sensodat-45s}

Dữ liệu là \textbf{SensoDat} (Birchler 2024, MSR --- ref {[}1{]}), kịch
bản SDC sinh tự động trên BeamNG, công khai. Ba split: \textbf{Train}
28.804 (FAIL 38,4\%), \textbf{Test} 7.202 (cùng phân phối),
\textbf{Competition} 956 --- đây là tập \textbf{OOD}: FAIL 36,9\%, đường
ngắn hơn hẳn (129--229m so với tới 454m ở train). Biểu đồ phải cho thấy
đường càng dài càng dễ FAIL. Khoảng độ dài hẹp hơn ở Competition chính
là \textbf{dịch phân phối} mà ta phải tổng quát hóa qua.

\begin{quote}
\textbf{VERIFY:} mọi số khớp \texttt{best.md}: 11059/28804=38,4\% ✓,
353/956=36,9\% ✓, 2765/7202=38,4\% ✓. Các thanh fail-rate theo độ dài là
nhóm tự tính (không có trong tracker) --- đúng xu hướng ``dài→FAIL
nhiều'', nhưng nếu thầy hỏi giá trị từng bin thì nói ``nhóm tính trực
tiếp trên split''.
\end{quote}

\subsection{s07 · Ví dụ kịch bản ---
{[}35s{]}}\label{s07-vuxed-dux1ee5-kux1ecbch-bux1ea3n-35s}

Đây là dữ liệu thật từ SensoDat và BeamNG. Quan sát: PASS thường là cong
mượt; FAIL thường có \textbf{chicane} --- cua zigzag, độ cong đổi dấu
nhanh. Histogram bên phải: FAIL lệch về phía ``tổng độ đổi hướng'' lớn
hơn. Điểm mấu chốt: \textbf{không có tọa độ (x,y) tuyệt đối nào quyết
định nhãn --- chỉ hình dạng}. Đây là gợi ý đầu tiên cho ràng buộc bất
biến.

\begin{quote}
\textbf{VERIFY:} ảnh thật (SensoDat + BeamNG.tech) --- đã có dẫn nguồn ở
slide Sources ✓.
\end{quote}

\subsection{s08 · Related Work --- 3 hướng ---
{[}35s{]}}\label{s08-related-work-3-hux1b0ux1edbng-35s}

Ba hướng tiếp cận trước đây: (1) \textbf{Search-based/diversity} (GA,
greedy) --- chọn test đa dạng, nhưng đa dạng ≠ lộ lỗi; (2)
\textbf{Feature-based ML} (ITS4SDC: bi-LSTM trên 2 đặc trưng chuỗi) ---
rẻ, dễ giải thích, nhưng ít đặc trưng và phụ thuộc khung; (3)
\textbf{Deep learning} (GNN, ResNet, RoadFury) --- APFD cao nhất nhưng
\textbf{không bất biến}, xoay/lấy mẫu là điểm trôi. Khoảng trống: chưa
method nào \textbf{đảm bảo} bất biến --- tất cả dựa vào augmentation
hoặc đặc trưng phụ thuộc khung.

\begin{quote}
\textbf{VERIFY:} ``RoadFury 0.804'' --- xem (!) 2 (số chuẩn là
0.807/0.8066). SO-SDC, SO-SDC ({[}2{]}, TOSEM'23) + Greedy đã verify
khớp literature. ITS4SDC là tool ICST'25 (arXiv 2501.03881, bi-LSTM)
nhưng \textbf{APFD 0.781 là nhóm tái lập} --- paper chỉ công bố F1=0.89.
Nên thêm 1 entry ref riêng cho ITS4SDC.
\end{quote}

\subsection{s09 · Related Work --- tổng hợp ---
{[}40s{]}}\label{s09-related-work-tux1ed5ng-hux1ee3p-40s}

Bảng tổng hợp APFD trên 956 test OOD. Đọc xu hướng: từ Random 0.493 lên
dần tới các deep model \textasciitilde0.80. Cột ``Bất biến'' thì
\textbf{mọi method trước đều ✗}. SE2RoadNet thêm một \textbf{trục đảm
bảo mới} (Δ=0), không chỉ chạy đua APFD. Đây là thông điệp: ta không cố
tăng APFD thêm vài phần nghìn, mà thêm một chiều bảo chứng chưa ai có.

\begin{quote}
\textbf{VERIFY:} (!) 2 (RoadFury 0.804 vs SE2RoadNet 0.805) + (!) 4
(nguồn các baseline). Câu ``thêm trục đảm bảo, không chỉ tăng APFD'' là
cách diễn đạt TRUNG THỰC nhất --- nên nhấn câu này thay vì khoe APFD cao
hơn.
\end{quote}

\begin{center}\rule{0.5\linewidth}{0.5pt}\end{center}

\section{§02 --- ROADFURY (Method phần 1, \textasciitilde2
phút)}\label{roadfury-method-phux1ea7n-1-2-phuxfat}

\subsection{s10 · Divider §02 --- {[}5s{]}}\label{s10-divider-02-5s}

Phần hai: RoadFury --- phương pháp nền tảng của nhóm.

\subsection{s11 · RoadFury pipeline ---
{[}30s{]}}\label{s11-roadfury-pipeline-30s}

RoadFury là sơ đồ TikZ này: chuẩn hóa đường về 197 điểm và \textbf{10
kênh đặc trưng} → RoadTransformer → SWA → inference (chấm điểm, sắp giảm
dần). Đây là baseline rất mạnh, nhưng có 2 điểm mù mà em phân tích ở
slide sau.

\subsection{s12 · RoadFury chi tiết ---
{[}55s{]}}\label{s12-roadfury-chi-tiux1ebft-55s}

Bên trái là 10 kênh đặc trưng: chiều dài đoạn, biến thiên góc, độ cong
Menger, jerk độ cong\ldots{} nhưng \textbf{3 kênh f5--f7 (sinθ, cosθ, θ
tuyệt đối) phụ thuộc khung quy chiếu}. Bên phải: RoadTransformer 829K
tham số --- Linear 10→128, token {[}CLS{]} + \textbf{PE vị trí tuyệt
đối}, 4 lớp Transformer, pool → sigmoid. Kết quả: APFD ≈ 0.804--0.807,
AUC 0.917, tốt nhất competition. Nhưng 3 kênh phụ thuộc khung + PE tuyệt
đối = \textbf{2 điểm mù}: APFD cao mà không có bảo chứng bất biến.

\begin{quote}
\textbf{VERIFY:} 829K params ✓ (828,801). AUC 0.917 ✓. APFD: deck ghi
``0.804 ± 0.012'' --- số canonical của dự án là \textbf{0.8066 ± 0.0124}
(best-single γ=2.5). Nói ``khoảng 0.805--0.807'' an toàn hơn ``0.804''.
Xem (!) 2.
\end{quote}

\subsection{s13 · Bridge ``Vá 2 điểm mù'' ---
{[}30s{]}}\label{s13-bridge-vuxe1-2-ux111iux1ec3m-muxf9-30s}

Mỗi điểm mù của RoadFury → một thành phần của SE2RoadNet. Điểm mù 1
(nhạy xoay, ITS4SDC Δ=0.057) → vá bằng \textbf{7 kênh bất biến SE(2)}
(bỏ sinθ, cosθ, θ). Điểm mù 2 (nhạy tần số lấy mẫu) → vá bằng
\textbf{attention bias theo Δs}. Tinh thần: chuyển từ
\textbf{augmentation (xấp xỉ)} sang \textbf{ràng buộc kiến trúc (bảo
chứng)} --- vá có chủ đích, không refactor mù.

\begin{quote}
\textbf{VERIFY:} Δ=0.057 cho ITS4SDC --- xem (!) 3 (xác nhận đã đo
thật). ``N=64 vs 197 → điểm đổi'' ✓ khớp tinh thần resolution probe.
\end{quote}

\begin{center}\rule{0.5\linewidth}{0.5pt}\end{center}

\section{§03 --- SE2ROADNET (Method phần 2 --- TRỌNG TÂM,
\textasciitilde4
phút)}\label{se2roadnet-method-phux1ea7n-2-trux1ecdng-tuxe2m-4-phuxfat}

\subsection{s14 · Divider §03 --- {[}5s{]}}\label{s14-divider-03-5s}

Phần ba: SE2RoadNet --- đề xuất chính.

\subsection{s15 · SE2RoadNet overview ---
{[}30s{]}}\label{s15-se2roadnet-overview-30s}

Sơ đồ tổng: 7 kênh bất biến → \textbf{InvariantBlock × 6} (attention
bias theo Δs) → pool → Failure Score. Kết quả then chốt ngay đây:
\textbf{rotation probe ΔAPFD = 0.0000} --- xoay đường thì điểm không đổi
đến từng bit. Em đi vào 4 bước.

\begin{quote}
\textbf{VERIFY:} Δ=0.0000 ✓ (tracker Exp 02). Caption không nhắc
resolution ở đây --- tốt.
\end{quote}

\subsection{s16 · Ý tưởng cốt lõi SE(2) ---
{[}45s{]}}\label{s16-uxfd-tux1b0ux1edfng-cux1ed1t-luxf5i-se2-45s}

Định lý mục tiêu: f\_θ(R·R + t) = f\_θ(R) \textbf{đến từng bit}, với mọi
phép xoay R và tịnh tiến t. Bốn panel minh họa cùng một con đường ---
gốc, xoay, tịnh tiến, đổi gốc tọa độ --- đều cho \textbf{Score =
0.8047}. Bất biến ở đây là \textbf{tính chất của kiến trúc}, không phải
của data augmentation. Đó là khác biệt cơ bản: augmentation chỉ làm mô
hình ``quen'' với phép xoay; kiến trúc của ta \textbf{không cho phép}
điểm số đổi.

\begin{quote}
\textbf{VERIFY:} 0.8047 = single-pass thật ✓. ``bit-identical'' đúng về
APFD ở precision báo cáo; ở mức logit float32 có sai số
\textasciitilde1e-7 do ma trận xoay R chứa sin/cos (xem Q\&A Q6 backup).
Nếu thầy bắt bẻ ``bit'', dùng câu trả lời Q6.
\end{quote}

\subsection{s17 · Bước 1 --- 7 kênh bất biến ---
{[}55s{]}}\label{s17-bux1b0ux1edbc-1-7-kuxeanh-bux1ea5t-biux1ebfn-55s}

Bước 1: chỉ giữ \textbf{đại lượng nội tại} của đường: (1) Δs chiều dài
đoạn, (2) \textbar Δθ\textbar{} biến thiên góc, (3) κ độ cong, (4) dκ/ds
jerk, (5) d²κ/ds², (6) s\_norm arclength chuẩn hóa, (7) σ\_local độ lệch
chuẩn cục bộ. \textbf{0 kênh phụ thuộc khung} --- đã bỏ sinθ, cosθ, θ.
Cơ sở lý thuyết là \textbf{Frenet--Serret}: một đường cong phẳng được
xác định duy nhất đến phép xoay-tịnh tiến bởi hàm độ cong κ(s). Nên bỏ
(x,y) và hướng tuyệt đối \textbf{không mất thông tin} --- chỉ bỏ phần dư
thừa gây nhạy xoay. Bất biến được đảm bảo \textbf{ngay từ khâu đặc
trưng}, trước khi mô hình học.

\begin{quote}
\textbf{VERIFY:} 7 kênh khớp tracker (``7-channel SE(2)-invariant
features'') ✓.
\end{quote}

\subsection{s18 · Bước 2 --- kiến trúc InvariantBlock×6 ---
{[}45s{]}}\label{s18-bux1b0ux1edbc-2-kiux1ebfn-truxfac-invariantblock6-45s}

Bước 2: kiến trúc. Mỗi điểm → token 192 chiều. Thêm token {[}CLS{]}.
\textbf{6 InvariantBlock}, mỗi block = MHA 8 heads + FFN 512 +
LayerNorm, kèm \textbf{relative-arclength bias} (bước 3). Pool {[}CLS{]}
→ head 192→64→1 → xác suất FAIL. Tổng \textbf{2.11M tham số}
(\textasciitilde2.5× RoadFury) --- vẫn nhẹ, train 24 phút. Một config
duy nhất, không tinh chỉnh theo từng dataset.

\begin{quote}
\textbf{VERIFY:} 2.11M ✓ (2,108,721), 24.2 phút ✓, d=192/depth6/8heads
✓.
\end{quote}

\subsection{s19 · Bước 3 --- Attention bias Δs --- {[}60s{]} {[}*{]}
phần
lõi}\label{s19-bux1b0ux1edbc-3-attention-bias-ux3b4s-60s-phux1ea7n-luxf5i}

Bước 3 là lõi kỹ thuật. \textbf{Vấn đề của PE tuyệt đối}: PE chuẩn dùng
sin/cos theo \emph{index} của token, khóa mô hình vào vị trí tuyệt đối;
lấy mẫu khác thì index ứng arclength khác → vỡ bất biến. \textbf{Giải
pháp}: thêm bias vào ma trận attention chỉ phụ thuộc \textbf{hiệu số
arclength} Δs\_ij = s\_i − s\_j, qua MLP(sin(Δs·ω)) với ω là 32 tần số
Fourier ngẫu nhiên cố định (RFF). Vì Δs là đại lượng nội tại nên bias
\textbf{bất biến} xoay/tịnh tiến. Heatmap: attention giảm theo
\textbar Δs\textbar, không bao giờ phụ thuộc i, j riêng lẻ. Giá phải
trả: chi phí O(B·L²·32) mỗi block --- đây là điểm tốn nhất, nên train 24
phút; hướng cải tiến là \textbf{RoPE-1D} trên s\_norm.

\begin{quote}
\textbf{VERIFY:} O(B·L²·32) ✓, RoPE-1D là hướng tương lai ✓ (tracker
action item).
\end{quote}

\subsection{s20 · Bước 4 --- Huấn luyện ---
{[}45s{]}}\label{s20-bux1b0ux1edbc-4-huux1ea5n-luyux1ec7n-45s}

Bước 4: huấn luyện. \textbf{Focal loss γ=1.0}: vì FAIL chỉ
\textasciitilde30--40\%, BCE thuần khiến mô hình ``đoán PASS hết''; hệ
số (1−p̂\_t)\^{}γ giảm trọng số ca dễ, tăng trọng số ca khó. Cấu hình:
AdamW, LR 5e-4 cosine + warmup, batch 384, 80 epoch (SWA từ 56), bf16,
WeightedRandomSampler. Train 24,2 phút, 2.11M tham số. Sweep
γ∈{[}1,5{]}: APFD gần như phẳng (mọi γ đều trong ±σ) ⇒ chọn γ=1.0
(0.8095±0.0121). Không đổi siêu tham số giữa các bench.

\begin{quote}
\textbf{VERIFY:} γ=1.0 (deck đã đổi theo sweep γ, 2026-07-08;
s18/s23/s24 vẫn hiển 0.8047/0.805/0.934 từ lần chạy γ=1.5 --- cascade
chưa đồng bộ), SWA từ ep56 ✓, batch 384 ✓, 24.2 phút ✓. Lưu ý: RoadFury
best-single dùng γ=2.5, SE2RoadNet nay dùng γ=1.0 --- cả hai đều đúng
trong ngữ cảnh riêng, đừng nhầm lẫn nếu thầy hỏi.
\end{quote}

\begin{center}\rule{0.5\linewidth}{0.5pt}\end{center}

\section{§04 --- ĐÁNH GIÁ (Benchmark --- TRỌNG TÂM, \textasciitilde3.5
phút)}\label{ux111uxe1nh-giuxe1-benchmark-trux1ecdng-tuxe2m-3.5-phuxfat}

\subsection{s21 · Divider §04 --- {[}5s{]}}\label{s21-divider-04-5s}

Phần bốn: đánh giá thực nghiệm.

\subsection{s22 · APFD \& giao thức ---
{[}50s{]}}\label{s22-apfd-giao-thux1ee9c-50s}

\textbf{APFD} = 1 − (ΣTF\_i)/(n·m) + 1/(2n), khoảng {[}0,1{]}, càng cao
càng tốt; ngẫu nhiên \textasciitilde0.5. Ta đánh giá 3 lớp: (1)
\textbf{single-pass} trên cả 956 test → APFD = 0.8047; (2)
\textbf{multi-trial 30 lần}, mỗi lần lấy ngẫu nhiên 287/956 → 0.8048 ±
0.0118, để loại ``may rủi'' do thứ tự cố định; (3) \textbf{rotation
probe} --- xoay cả split bằng 6 góc rồi đo Δ = max − min APFD, để kiểm
chứng bất biến.

\begin{quote}
\textbf{VERIFY:} công thức ✓, 0.8047 ✓, 0.8048±0.0118 ✓, 287/956 ✓
(=max(50, 0.3·956)). Tất cả khớp tracker.
\end{quote}

\subsection{s23 · Headline Δ=0 --- {[}70s{]} {[}!!{]} slide quan trọng
nhất}\label{s23-headline-ux3b40-70s-slide-quan-trux1ecdng-nhux1ea5t}

Đây là kết quả headline. Bảng: 6 góc xoay (0, +30, +60, +90, +180, −45),
SE2RoadNet đều cho \textbf{0.8047} --- Δ = \textbf{0.0000}. Để so sánh,
ITS4SDC tụt từ 0.781 xuống tận 0.724 ở +30°, Δ = 0.057. Em nhấn: đây
\textbf{không phải ``trong sai số 1e-6''}, mà bằng nhau \textbf{đến từng
bit float} --- vì pipeline 7 kênh trả vector y hệt sau xoay, forward
pass deterministic, nên logit y nguyên. Đây là một trong số rất ít chỗ
trong ML có thể nói ``lý thuyết được xác minh bằng thực nghiệm'' mà
không cần đính kèm sigma.

\begin{quote}
\textbf{VERIFY:} Δ=0.0000 ✓✓ (tracker, kết quả sạch nhất dự án).
\textbf{Bảng ITS4SDC 6 góc (0.781→0.724\ldots) --- xem (!) 3}: xác nhận
đã đo thật, không phải số minh họa trong dải dự đoán 0.04--0.07.
\end{quote}

\subsection{s24 · Leaderboard --- {[}60s{]}}\label{s24-leaderboard-60s}

Bảng xếp hạng APFD với 8 baseline. SE2RoadNet 0.805 --- \textbf{ngang
baseline tốt nhất trong sai số}, nhưng: AUC tăng từ 0.917 lên 0.934,
\textbf{và} thêm bảo chứng Δ=0 mà không method nào có → cải thiện theo
nghĩa \textbf{Pareto}: không tệ hơn ở chiều nào, tốt hơn ở AUC và lý
thuyết. Không đánh đổi: giữ APFD đỉnh, tăng AUC, là phương pháp duy nhất
bất biến SE(2) tuyệt đối.

\begin{quote}
\textbf{VERIFY:} (!) 2 --- RoadFury phải là \textbf{0.807} (không
0.804), nên nói ``ngang trong sai số'' chứ KHÔNG nói ``cao hơn''. AUC
0.934 → thực ra 0.9347≈\textbf{0.935}. ``+0.017'' ✓ (0.9347−0.9170). Đây
là slide thầy dễ soi nhất → bám đúng từ ``Pareto / ngang trong σ'', đừng
nói ``thắng APFD''.
\end{quote}

\subsection{s25 · AUC vs APFD --- {[}45s{]}}\label{s25-auc-vs-apfd-45s}

Vì sao tách bạch AUC và APFD? \textbf{AUC} đo khả năng phân loại
PASS/FAIL trên từng cặp; \textbf{APFD} đo tốc độ lộ lỗi theo thứ tự chạy
--- đúng cái Competition tối ưu. Trong gần như mọi thí nghiệm của nhóm,
hai chỉ số này \textbf{phân kỳ}: SE2RoadNet tăng AUC 0.917→0.934 nhưng
APFD gần như phẳng \textasciitilde0.805. Bài học: khi báo cáo phải nói
rõ \textbf{metric nào} đang tối ưu, không gộp thành ``một điểm tốt
hơn''.

\begin{quote}
\textbf{VERIFY:} Phân kỳ AUC/APFD ✓ (tracker xác nhận ở Exp
01,02,03,04). Đây là phần TRUNG THỰC giúp giải thích vì sao APFD không
tăng --- nên giữ và nói rõ.
\end{quote}

\begin{center}\rule{0.5\linewidth}{0.5pt}\end{center}

\section{§05 --- HẠN CHẾ \& HƯỚNG PHÁT TRIỂN (\textasciitilde1.5
phút)}\label{hux1ea1n-chux1ebf-hux1b0ux1edbng-phuxe1t-triux1ec3n-1.5-phuxfat}

\subsection{s26 · Divider §05 --- {[}5s{]}}\label{s26-divider-05-5s}

Phần năm: hạn chế và hướng phát triển.

\subsection{s27 · Hạn chế ---
{[}50s{]}}\label{s27-hux1ea1n-chux1ebf-50s}

Em xin thẳng thắn. \textbf{Về chỉ số}: AUC và APFD phân kỳ; listwise
loss chưa tăng APFD trung bình (chỉ giảm σ → đóng góp ``ổn định''); ở
bench FAIL cao (vd 95\% FAIL) APFD \textasciitilde0.52 là \textbf{trần},
không phải thất bại. \textbf{Về phương pháp}: conformal an toàn còn dở
(v1 valid-nhưng-vô nghĩa, v2 informative-nhưng-invalid, cần v3);
IRM/TENT chưa đóng được gap SensoDat→Competition (negative đã biết); bất
biến lấy mẫu mới là \textbf{xấp xỉ} (Δ≈0.0012), chưa exact như phép
xoay. Tóm lại: bất biến xoay là exact, các đảm bảo khác đang hoàn thiện
--- báo cáo trung thực thay vì giấu.

\begin{quote}
\textbf{VERIFY:} TẤT CẢ khớp tracker ✓ (listwise σ↓, conformal v1/v2,
IRM/TENT negative, ceiling 0.52). Resolution Δ≈0.0012 --- xem (!) 5 (đó
là số FNO).
\end{quote}

\subsection{s28 · Hướng phát triển ---
{[}40s{]}}\label{s28-hux1b0ux1edbng-phuxe1t-triux1ec3n-40s}

Bốn hướng: (1) tổng quát hóa \textbf{8+ benchmark} (OOB, SDC-Scissor,
sdc-travel) với một công thức, không tinh chỉnh; (2) \textbf{conformal
v3} vừa valid vừa non-vacuous; (3) \textbf{bất biến lấy mẫu exact} bằng
RoPE-1D thay RFF; (4) \textbf{SSL vật lý} (physics-informed) thay SSL
hình học ngây thơ (đã chuyển kém). Mục tiêu: từ một baseline bất biến
cho SDC → công thức chung bất biến và audit-readable.

\begin{quote}
\textbf{VERIFY:} ``8+ benchmark'' đang là \textbf{hướng tương lai},
không phải kết quả đã có --- slide ghi đúng là future work ✓. Đừng lỡ
miệng nói ``đã chạy 8 bench''.
\end{quote}

\subsection{s29 · Tham khảo --- {[}15s{]}}\label{s29-tham-khux1ea3o-15s}

Tài liệu tham khảo theo IEEE. Nguồn dữ liệu là SensoDat {[}1{]}; SO-SDC
{[}2{]}, giao thức APFD {[}3{]}{[}4{]}; nền tảng kỹ thuật: Attention
{[}5{]}, Focal {[}6{]}, SWA {[}7{]}, equivariance {[}8{]}, RFF {[}9{]},
conformal {[}10{]}. RoadFury và SE2RoadNet là của nhóm.

\subsection{s30 · Video --- {[}20s{]}}\label{s30-video-20s}

(Track Method 60\%) Đây là video demo trình bày phương pháp; link Google
Drive và Facebook ở đây. \emph{(Bấm vào nếu cần phát.)}

\begin{quote}
\textbf{VERIFY:} thumbnail vẫn là placeholder --- \textbf{thay ảnh thật
trước khi nộp}. Link đã nhúng sẵn ✓.
\end{quote}

\subsection{s31 · Cảm ơn --- {[}10s{]}}\label{s31-cux1ea3m-ux1a1n-10s}

Em cảm ơn thầy và cả lớp đã lắng nghe. Em xin nhận câu hỏi.

\begin{center}\rule{0.5\linewidth}{0.5pt}\end{center}

\section{BACKUP (s32--s36) --- chỉ mở khi Q\&A
hỏi}\label{backup-s32s36-chux1ec9-mux1edf-khi-qa-hux1ecfi}

\subsection{s32 · B1 --- APFD ví dụ
tay}\label{s32-b1-apfd-vuxed-dux1ee5-tay}

n=5, m=2 FAIL. Xếp tốt (FAIL ở 1,2): APFD = 1 − 3/10 + 1/10 =
\textbf{0.80}. Xếp kém (FAIL ở 4,5): 1 − 9/10 + 1/10 = \textbf{0.20}.
Cùng tập lỗi, chỉ khác thứ tự, APFD chênh 4 lần --- đó là toàn bộ giá
trị của test prioritization.

\begin{quote}
\textbf{VERIFY:} số học ✓ đúng tuyệt đối.
\end{quote}

\subsection{s33 · B2 --- Multi-trial}\label{s33-b2-multi-trial}

Vì sao 30 lần: single-pass nhạy với thứ tự cố định → dễ may rủi. Mỗi
trial lấy ngẫu nhiên max(50, 0.3·\textbar test\textbar) = 287/956, tính
APFD độc lập; báo cáo mean ± σ qua 30 trials, seed 42. σ thường quan
trọng hơn mean (độ ổn định). SE2RoadNet = 0.8048 ± 0.0118.

\begin{quote}
\textbf{VERIFY:} mọi số ✓ khớp code \texttt{multi\_trial} trong best.md.
\end{quote}

\subsection{s34 · B3 --- Resolution
probe}\label{s34-b3-resolution-probe}

Resample cùng đường ở N ∈ \{64..197\}, đo Δ = max − min APFD. Kết quả Δ
≈ 0.0012 --- rất nhỏ nhưng \textbf{chưa exact} như phép xoay (Δ=0). Lý
do: RFF bias theo Δs là gần-bất-biến tại tham số hóa → hướng RoPE-1D.

\begin{quote}
\textbf{VERIFY:} (!) 5 --- Δ=0.0012 đo trên \textbf{FNO (Exp 01)}, và
bar (.8060/.8066/ .8072/.8067) KHÔNG khớp bảng FNO thật
(.8051/.8060/.8063/.8060/.8062). Nếu thầy hỏi: nói rõ ``probe resolution
chạy trên backbone FNO Exp 01; SE2RoadNet hiện chỉ có exact rotation,
resolution đang là hướng''.
\end{quote}

\subsection{s35 · B4 --- Chi phí quy mô
lớn}\label{s35-b4-chi-phuxed-quy-muxf4-lux1edbn}

Inference chỉ 1 forward pass/đường → chấm cả tập trong vài giây GPU. Chi
phí thật nằm ở \textbf{mô phỏng vật lý} của các test được chọn, không
phải scorer. Train một lần 24,2 phút, 2.11M tham số → tái lập rẻ.

\begin{quote}
\textbf{VERIFY:} điểm 3 (``cắt 50--80\% số kịch bản'') --- \textbf{cùng
vấn đề (!) 1}. 24.2 phút ✓, 2.11M ✓.
\end{quote}

\subsection{s36 · B5 --- Chứng minh bất
biến}\label{s36-b5-chux1ee9ng-minh-bux1ea5t-biux1ebfn}

SE(2): p → R·p + t. Δs bất biến (R trực giao bảo toàn khoảng cách, t
triệt tiêu); \textbar Δθ\textbar{} bất biến (R bảo toàn góc); κ và đạo
hàm theo s là hàm của đại lượng bất biến → bất biến. Bỏ sinθ, cosθ, θ là
điều kiện cần. Hệ quả: F(R·R+t) = F(R) pixel-identical → mọi tầng sau
chỉ nhận đầu vào bất biến → output bất biến.

\begin{quote}
\textbf{VERIFY:} lập luận toán đúng ✓. Đây là backup mạnh cho Q1.
\end{quote}

\begin{center}\rule{0.5\linewidth}{0.5pt}\end{center}

\section{CHUẨN BỊ Q\&A (mang sẵn trong
đầu)}\label{chuux1ea9n-bux1ecb-qa-mang-sux1eb5n-trong-ux111ux1ea7u}

\textbf{Q: ``Cắt 50--80\% thời gian --- số này ở đâu?''} → Trung thực:
đây là khoảng quen thuộc của lĩnh vực regression/test prioritization,
không phải số nhóm đo trực tiếp. Số NHÓM đo được: với ranking của
SE2RoadNet, prefix K=50/287 (\textasciitilde17\% test) đã đạt APFD@K ≈
0.94 → phần lớn giá trị phát hiện lỗi nằm ở một phần nhỏ test đầu.
\emph{(Nên sửa slide theo (!) 1 để khỏi bị hỏi.)}

\textbf{Q: ``SE2RoadNet APFD thấp hơn baseline 0.0018, sao gọi là tốt
hơn?''} → Em \textbf{không} claim APFD tốt hơn. Em claim \textbf{Pareto
improvement}: APFD ngang trong sai số (σ≈0.012), AUC cao hơn rõ rệt
(0.917→0.934), \textbf{cộng} guarantee Δ=0 mà không method nào có. Trong
bối cảnh safety-critical, guarantee đáng giá hơn 0.0018 APFD.

\textbf{Q: ``Vì sao Δ=0 gọi là exact --- mạng nơ-ron mà?''} → Bất biến
đến từ \textbf{pipeline đặc trưng}, không từ tham số học được. 7 kênh
không chứa (x,y) hay sin/cos tuyệt đối, nên xoay đường rồi trích lại cho
vector y hệt; mạng đã eval (dropout tắt), deterministic → logit y hệt →
APFD y hệt. Ở mức logit float32 vẫn có \textasciitilde1e-7 do ma trận
xoay R chứa sin/cos không exact, nhưng APFD ở precision báo cáo là
bit-identical.

\textbf{Q: ``O(L²) với L=197 thì scale thế nào?''} → Đang triển khai
\textbf{RoPE-1D} trên s\_norm thay RFF bias: giữ bất biến arclength mà
chi phí O(L), kỳ vọng giảm train từ 24 xuống 3--4 phút.

\textbf{Q: ``Competition sao gọi là OOD?''} → (1) generator khác; (2)
FAIL rate khác (36,9\% vs 38,4\%); (3) độ dài khác hẳn (129--229m vs tới
454m). Bằng chứng: APFD SensoDat-test \textasciitilde0.76 vs Competition
\textasciitilde0.80, chênh \textgreater0.04 → dịch phân phối rõ.

\textbf{Q: ``Các baseline 0.487/0.533/\ldots{} lấy ở đâu?''} → SO-SDC
0.765 + Greedy 0.795 khớp chính xác Birchler TOSEM'23 ({[}2{]}). ITS4SDC
là tool ICST'25 (bi-LSTM, arXiv 2501.03881) nhưng APFD 0.781 là nhóm tái
lập (paper chỉ có F1=0.89). Random/LLM/GNN/ResNet là nhóm tái lập trong
cùng harness (956 test OOD, 30 trials). \emph{(Chuẩn bị sẵn log tái lập
--- xem (!) 4.)}

\end{document}
